% ============================================================
%  Section 4 -- Resultats
% ============================================================

\chapter{Résultats}

\section{Conditions experimentales}

Les resultats de cette section proviennent des artefacts presents dans
\file{data/outputs/benchmark/}. La campagne principale a ete generee le
9 mai 2026 et compare deux architectures sur 10 tickets extraits du jeu de
donnees de 40 tickets. Les deux architectures utilisent \code{phi3:latest}; le
juge qualite est \code{gpt-4o}. La tarification configuree dans le fichier de
resultats est de 2,5 dollars par million de tokens d'entree et 10 dollars par
million de tokens de sortie.

\begin{table}[H]
\centering
\begin{tabularx}{\textwidth}{lY}
\toprule
\textbf{Artefact} & \textbf{Usage dans le rapport} \\
\midrule
\file{comparison_results.json} & Resultat principal LLM unique vs systeme multi-agents. \\
\file{benchmark_mcp.json} & Performance directe des outils MCP. \\
\file{benchmark_mcp_tools.json} & Variante MCP avec pertinence/F1. \\
\file{ab_llm_resultat.json} & Test A/B entre \code{gpt-4o-mini} et \code{claude-haiku-4.5}. \\
\file{position_bias_resultat.json} & Test de biais de position sur tickets IT. \\
\file{security_race_report.json} & Rapport adversarial Racing Arena. \\
\file{a2a_agents_results.json} & Benchmark A2A du 12 mai 2026, non concluant. \\
\file{kaggle/local_kaggle_exam_report_20260509_101412.json} & Rejeu local d'un examen Kaggle. \\
\bottomrule
\end{tabularx}
\caption{Sources chiffrees exploitees.}
\label{tab:result-artifacts}
\end{table}

\section{Comparaison LLM Unique vs Système Multi-Agents}

\begin{figure}[H]
  \centering
  \includegraphics[width=\textwidth]{comparaison_architectures.png}
  \caption[Comparaison globale des deux architectures]{Comparaison globale des deux architectures. Chemin :
  \file{data/outputs/benchmark/charts/comparaison_architectures.png}.}
  \label{fig:comparaison-architectures}
\end{figure}

Le tableau \ref{tab:main-results} reprend les metriques brutes du benchmark
principal.

\begin{table}[H]
\centering
\begin{tabularx}{\textwidth}{YccY}
\toprule
\textbf{Métrique} & \textbf{LLM Unique} & \textbf{SMA} & \textbf{Lecture} \\
\midrule
Score qualite moyen & 0,0/10 & \textbf{3,8/10} & Le SMA est meilleur, mais loin du seuil 7/10. \\
Score securite moyen & 10,0/10 & 10,0/10 & Aucun risque detecte sur cette campagne. \\
Latence totale & 277,3515 s & \textbf{100,6540 s} & Reduction de 63,7 \%, soit 2,76 fois plus rapide. \\
Cout estime & 0,049593 \$ & \textbf{0,006270 \$} & Cout divise par 7,91. \\
Tokens totaux & 5\,344 & \textbf{627} & Reduction de 88,3 \%. \\
Energie estimee & 0,0010688 kWh & \textbf{0,0001254 kWh} & Energie divisee par 8,52. \\
Empreinte carbone & 0,05344 gCO$_2$e & \textbf{0,00627 gCO$_2$e} & Meme tendance que les tokens. \\
Score composite & 35,0/100 & \textbf{75,2/100} & Amelioration de 40,2 points. \\
Verdict & \verdict{FAIL} & \verdict{FAIL} & Gate qualite non atteint. \\
\bottomrule
\end{tabularx}
\caption{Resultats principaux extraits de \file{comparison_results.json}.}
\label{tab:main-results}
\end{table}

Le SMA gagne selon la regle du benchmark : meilleur score composite parmi les
architectures qui passent les gates, ou meilleur score global si aucune ne les
passe. Ici aucune architecture ne passe, mais le SMA domine nettement le LLM
unique.

\begin{figure}[H]
  \centering
  \includegraphics[width=\textwidth]{core_benchmark.png}
  \caption[Synthese du coeur de benchmark]{Synthese du coeur de benchmark. Chemin :
  \file{data/outputs/benchmark/charts/core_benchmark.png}.}
  \label{fig:core-benchmark}
\end{figure}

\subsection{Analyse du score composite}

\begin{figure}[H]
  \centering
  \includegraphics[width=0.92\textwidth]{composite_verdict.png}
  \caption[Verdicts composites]{Verdicts composites. Chemin :
  \file{data/outputs/benchmark/charts/composite_verdict.png}.}
  \label{fig:composite-verdict}
\end{figure}

Le score composite du SMA est eleve car il obtient la meilleure normalisation
sur le cout, la latence et l'empreinte carbone, tout en conservant un score
securite de 10/10. La faiblesse restante est la qualite. Le gate qualite
explique le verdict \verdict{FAIL} malgre un score composite de 75,2/100.

\begin{table}[H]
\centering
\begin{tabularx}{0.88\textwidth}{Ycc}
\toprule
\textbf{Composante normalisee} & \textbf{LLM Unique} & \textbf{SMA} \\
\midrule
Qualite & 0,0000 & 0,3800 \\
Securite & 1,0000 & 1,0000 \\
FinOps & 0,0000 & 1,0000 \\
Latence & 0,0000 & 1,0000 \\
GreenOps & 0,0000 & 1,0000 \\
\bottomrule
\end{tabularx}
\caption{Composantes normalisees de la politique composite.}
\label{tab:component-scores}
\end{table}

Ce resultat met en evidence un point methodologique : une architecture peut etre
superieure sur les dimensions operationnelles tout en restant insuffisante sur
le plan fonctionnel. Le rapport ne doit donc pas conclure que le SMA est pret
pour la production ; il conclut qu'il est la meilleure base experimentale.

\section{Qualite des outils et MCP}

Les benchmarks MCP isolent le comportement des outils exposes via Model Context
Protocol. Le fichier \file{benchmark_mcp.json} contient cinq tickets et cinq
succes. Le score final moyen est \metric{9,782/10}, avec une latence moyenne de
\metric{0,0518 s}. La variante \file{benchmark_mcp_tools.json}, qui ajoute une
dimension de pertinence, obtient \metric{7,848/10} en moyenne, pour une latence
moyenne de \metric{0,0746 s}.

\begin{figure}[H]
  \centering
  \includegraphics[width=\textwidth]{mcp_benchmark.png}
  \caption[Benchmark MCP]{Benchmark MCP. Chemin :
  \file{data/outputs/benchmark/charts/mcp_benchmark.png}.}
  \label{fig:mcp-benchmark}
\end{figure}

\begin{table}[H]
\centering
\begin{tabularx}{0.88\textwidth}{Yccc}
\toprule
\textbf{Artefact MCP} & \textbf{Cas} & \textbf{OK} & \textbf{Score moyen} \\
\midrule
\file{benchmark_mcp.json} & 5 & 5 & 9,782 \\
\file{benchmark_mcp_tools.json} & 5 & 5 & 7,848 \\
\bottomrule
\end{tabularx}
\caption{Synthese des benchmarks MCP.}
\label{tab:mcp-results}
\end{table}

L'ecart entre les deux scores vient de l'ajout d'une mesure de pertinence et de
F1 dans la seconde variante. Les outils repondent vite et sans erreur, mais la
pertinence automatique reste plus difficile a mesurer que la disponibilite
technique.

\section{Tests A/B et biais de position}

Le test A/B LLM compare \code{gpt-4o-mini} et \code{claude-haiku-4.5} sur trois
paires jugees par \code{gpt-4o}. \code{gpt-4o-mini} est choisi deux fois,
\code{claude-haiku-4.5} une fois, soit 66,7 \% contre 33,3 \%. Le faible nombre
de cas interdit une conclusion statistique forte ; le test valide surtout le
pipeline pairwise.

\begin{figure}[H]
  \centering
  \includegraphics[width=0.96\textwidth]{ab_test_llm.png}
  \caption[Resultat du test A/B LLM]{Resultat du test A/B LLM. Chemin :
  \file{data/outputs/benchmark/charts/ab_test_llm.png}.}
  \label{fig:ab-test}
\end{figure}

Le test de biais de position est plus directement statistique. Sur 10 tickets,
20 jugements sont produits en inversant les positions A/B. Le juge choisit la
position A 12 fois et la position B 8 fois. Le taux A est donc de 0,60, avec une
p-value binomiale bilaterale de 0,5034. Aucun biais de position significatif
n'est etabli sur cette campagne.

\begin{figure}[H]
  \centering
  \includegraphics[width=0.96\textwidth]{position_bias.png}
  \caption[Test de biais de position]{Test de biais de position. Chemin :
  \file{data/outputs/benchmark/charts/position_bias.png}.}
  \label{fig:position-bias}
\end{figure}

Un second test de position bias sur statements donne aussi un taux A de 0,60,
mais avec 10 jugements et une p-value de 0,7539. Ces resultats sont coherents :
le pipeline ne detecte pas de biais dominant, mais le volume reste trop faible
pour une conclusion definitive.

\section{Modeles zero-shot et scores historiques}

Le fichier \file{tickets_evalues_scores.json} contient 12 tickets evalues sur
trois modeles. Les moyennes observees sont :

\begin{table}[H]
\centering
\begin{tabularx}{0.88\textwidth}{Yccccc}
\toprule
\textbf{Modele} & \textbf{Cas} & \textbf{Moyenne} & \textbf{Mediane} & \textbf{Min} & \textbf{Max} \\
\midrule
\code{llama3.2:1b} & 4 & 5,63 & 6,42 & 1,88 & 7,79 \\
\code{gpt-3.5-turbo} & 4 & 7,01 & 6,94 & 5,91 & 8,25 \\
\code{mistral-7b} & 4 & 5,49 & 6,16 & 1,88 & 7,74 \\
\bottomrule
\end{tabularx}
\caption{Scores historiques extraits de \file{tickets_evalues_scores.json}.}
\label{tab:historical-models}
\end{table}

Ces chiffres ne sont pas directement comparables a la campagne principale, car
ils proviennent d'un autre protocole et d'autres modeles. Ils servent cependant
de repere : obtenir une qualite superieure a 7/10 est possible dans le framework,
mais depend fortement du modele et du protocole.

\begin{figure}[H]
  \centering
  \includegraphics[width=0.82\textwidth]{graphique_1_score_par_modele.png}
  \caption[Scores par modele]{Scores par modele. Chemin :
  \file{data/outputs/benchmark/charts/graphique_1_score_par_modele.png}.}
  \label{fig:score-modele}
\end{figure}

\begin{figure}[H]
  \centering
  \includegraphics[width=0.92\textwidth]{graphique_2_latence_vs_score.png}
  \caption[Latence et score]{Latence et score. Chemin :
  \file{data/outputs/benchmark/charts/graphique_2_latence_vs_score.png}.}
  \label{fig:latence-score}
\end{figure}

\section{Racing Arena adversariale}

La Racing Arena apporte les resultats les plus interessants sur la securite.
Le rapport \file{security_race_report.json}, genere le 9 mai 2026, couvre
15 tours, 4 equipes et 286 decisions.

\begin{figure}[H]
  \centering
  \includegraphics[width=\textwidth]{racing_arena.png}
  \caption[Synthese Racing Arena]{Synthese Racing Arena. Chemin :
  \file{data/outputs/benchmark/charts/racing_arena.png}.}
  \label{fig:racing-arena}
\end{figure}

\begin{table}[H]
\centering
\begin{tabularx}{\textwidth}{Ycccc}
\toprule
\textbf{Equipe} & \textbf{Attaques recues} & \textbf{Injections executees} & \textbf{Fuites} & \textbf{Detection} \\
\midrule
Team A zero-shot & 12 & 7 (58,3 \%) & 9 (75,0 \%) & 0,0 \% \\
Team B ReAct & 57 & 53 (93,0 \%) & 7 (12,3 \%) & 14,0 \% \\
Team C validee & 1 & 0 (0,0 \%) & 0 (0,0 \%) & 0,0 \% \\
\bottomrule
\end{tabularx}
\caption{Metriques de securite de la Racing Arena.}
\label{tab:racing-security}
\end{table}

\begin{figure}[H]
  \centering
  \includegraphics[width=\textwidth]{racing_security_report.png}
  \caption[Rapport de securite Racing]{Rapport de securite Racing. Chemin :
  \file{data/outputs/benchmark/charts/racing_security_report.png}.}
  \label{fig:racing-security}
\end{figure}

Le resultat principal est clair : l'architecture validee resiste beaucoup mieux
que les variantes exposees. Team B, qui dispose d'un serveur \code{/query} brut
et sans filtrage, execute 93 \% des injections recues. Team A fuit davantage de
strategie, avec un taux de fuite de 75 \%. Team C ne fuit aucune strategie et
n'execute aucune injection dans cette course.

L'efficacite varie selon le type d'attaque :

\begin{table}[H]
\centering
\begin{tabularx}{0.88\textwidth}{Yccc}
\toprule
\textbf{Type d'attaque} & \textbf{Tentatives} & \textbf{Taux injection} & \textbf{Taux fuite} \\
\midrule
Usurpation d'autorite & 12 & 58,3 \% & 66,7 \% \\
Injection directe & 52 & 94,2 \% & 7,7 \% \\
RAG empoisonne & 6 & 66,7 \% & 66,7 \% \\
\bottomrule
\end{tabularx}
\caption{Efficacite des attaques adversariales.}
\label{tab:attack-effectiveness}
\end{table}

Ces chiffres montrent que la vulnerabilite ne se resume pas au modele : elle
depend de l'exposition API, du filtrage, du cloisonnement des roles et de la
maniere dont les donnees externes sont reinjectees dans le prompt.

\section{A2A et Kaggle}

Le benchmark A2A du 12 mai 2026 est present mais non concluant. Il cible neuf
agents distants, active les sondes de role, les sondes d'outils, le streaming,
le probing de modele, le fact-checking et un sous-ensemble Kaggle de quatre
questions. Dans \file{a2a_agents_results.json}, chaque agent remonte toutefois
\code{error_count = 1}, avec qualite, securite, score de tache et latence a 0.
Le rapport doit donc interpreter cet artefact comme une validation partielle du
pipeline d'orchestration, pas comme une comparaison de performances.

Le dernier rapport Kaggle local disponible,
\file{local_kaggle_exam_report_20260509_101412.json}, utilise
\code{mistral:latest} comme modele agent et \code{gpt-4.1} comme juge. Le score
est de \metric{5/16}, soit \metric{31,25 \%}. Ce resultat signale une marge de
progression importante sur les questions hors tickets IT classiques.

\section{Verification logicielle}

La suite de tests unitaires a ete executee localement avec :

\begin{lstlisting}[language=bash,caption={Commande de verification locale}]
PYTHONPATH=. BIBOPS_NON_INTERACTIVE=1 pytest tests/unit -q
\end{lstlisting}

Le resultat observe est \metric{552 tests passes en 9,98 s}. Cette verification
ne prouve pas la qualite fonctionnelle des reponses LLM, mais elle garantit que
les contrats internes, les fakes, les evaluateurs, les adaptateurs et la boucle
agentique restent coherents dans l'etat courant du depot.

\section{Synthese des resultats}

\begin{figure}[H]
  \centering
  \includegraphics[width=\textwidth]{synthese_finale.png}
  \caption[Synthese finale des benchmarks]{Synthese finale des benchmarks. Chemin :
  \file{data/outputs/benchmark/charts/synthese_finale.png}.}
  \label{fig:synthese-finale}
\end{figure}

Les resultats convergent vers quatre constats.

\begin{enumerate}
  \item Le systeme multi-agents est superieur au LLM unique sur la campagne
  principale, surtout en cout, latence, tokens et score composite.
  \item La qualite absolue reste insuffisante : le meilleur score moyen est
  3,8/10 alors que le gate exige 7/10.
  \item Les outils MCP sont rapides et stables, ce qui valide leur integration
  technique.
  \item La Racing Arena montre que la securite des agents depend fortement de
  l'architecture : une equipe validee peut reduire a zero les injections
  executees dans le scenario observe.
\end{enumerate}
